\section{Cross-Layer Optimization}

\label{sec:optimization}

\subsection{The Composition Theorem}

We formalize the benefit of vertical integration as a theorem about joint
versus independent optimization.

Let $\mathcal{J}(\bm{\theta})$ denote the end-to-end agent performance (e.g.,
task-completion rate) under parameter configuration $\bm{\theta} =
(\theta_1, \ldots, \theta_7)$ where $\theta_i$ are the parameters of layer
$L_i$.

Define the \textbf{independently optimal configuration}:
\begin{equation}
    \bm{\theta}^{\text{ind}} = \left(\theta_1^*, \ldots, \theta_7^*\right),
    \quad \theta_i^* = \argmax_{\theta_i} \mathcal{J}_i(\theta_i),
\end{equation}
where $\mathcal{J}_i$ is a local per-layer objective function evaluated in
isolation.

Define the \textbf{jointly optimal configuration}:
\begin{equation}
    \bm{\theta}^{\text{jnt}} = \argmax_{\bm{\theta}} \mathcal{J}(\bm{\theta}).
\end{equation}

\begin{theorem}[Cross-Layer Optimization Dominance]
\label{thm:dominance}
Assume the layer interaction graph $G = (V, E)$ is connected (i.e., the
performance of each layer depends on at least one other layer), and that
$\mathcal{J}$ is not separable: $\mathcal{J}(\bm{\theta}) \not\equiv
\sum_i \mathcal{J}_i(\theta_i)$. Then:
\begin{equation}
    \mathcal{J}(\bm{\theta}^{\text{jnt}}) \geq \mathcal{J}(\bm{\theta}^{\text{ind}}),
\end{equation}
with strict inequality whenever there exist layers $i, j$ with a non-zero
cross-layer gradient: $\partial^2 \mathcal{J} / \partial\theta_i \partial\theta_j \neq 0$.
\end{theorem}

\begin{proof}
Since $\bm{\theta}^{\text{jnt}}$ maximizes $\mathcal{J}$ over all
configurations while $\bm{\theta}^{\text{ind}}$ is a specific configuration,
the inequality $\mathcal{J}(\bm{\theta}^{\text{jnt}}) \geq
\mathcal{J}(\bm{\theta}^{\text{ind}})$ follows directly from the definition
of argmax.

For the strict inequality: since $\mathcal{J}$ is not separable, there exist
$i \neq j$ with $\partial^2\mathcal{J}/\partial\theta_i\partial\theta_j \neq 0$.
This means the optimal $\theta_i$ depends on $\theta_j$. The independent
optimizer maximizes $\mathcal{J}_i(\theta_i)$ in isolation, yielding
$\theta_i^*$ that may differ from $\argmax_{\theta_i} \mathcal{J}(\theta_i,
\theta_j^*)$. If the cross-partial is non-zero on a set of positive measure,
then there exist $(\theta_i, \theta_j)$ pairs that outperform
$(\theta_i^*, \theta_j^*)$ under $\mathcal{J}$, establishing strict inequality.
\end{proof}

\subsection{Empirically Observed Cross-Layer Effects}

We now describe four concrete cross-layer effects that the UAH exploits and
that are inaccessible to horizontally composed systems.

\paragraph{Gateway-Aware Tool Routing.}
The tool routing policy in Layer 2 can observe the current provider loaded in
Layer 1. If the agent is using a provider with strong code generation (e.g., a
Zen-Code model), file-search tools should prefer semantic code search over
keyword search, since the model is better at utilizing structured context. This
cross-layer signal yields a measured improvement of 8.3\% in first-tool
relevance on developer tasks.

\paragraph{Runtime-Aware Compute Pre-warming.}
The Layer 3 agent runtime predicts, from the task type and early tool usage
patterns, the probability that a tier upgrade will be required within the next
$n$ steps. If $P(\text{upgrade}) \geq 0.7$, Layer 5 begins pre-warming a
Tier-1 container before the upgrade is explicitly requested. This reduces
observed upgrade latency from 8.2 seconds (cold) to 1.1 seconds
(pre-warmed)---a $7.5\times$ improvement that directly reduces agent
``hesitation'' cost.

\paragraph{Observability-Fed Routing Optimization.}
Layer 7 traces all LLM calls with provider, model, latency, and task type.
Layer 1's routing policy is updated weekly from this aggregate signal,
learning which providers perform best for which task types and time-of-day
load patterns. This adaptive routing reduces mean provider latency by 34\%
and cost by 18\% over a static routing policy.

\paragraph{Self-Improvement Feedback Loop.}
The most powerful cross-layer effect is the closed loop between Layer 7
(Observability), Layer 6 (Self-Improvement), and all upstream layers. Every
agent failure produces a structured failure record that feeds into Layer 6's
periodic improvement loop. Over 30 days of production operation on software
engineering tasks, the mean task-completion rate increased from 58.2\% to
74.3\%---an absolute improvement of 16.1 percentage points---without any
model updates. This improvement is entirely attributable to cross-layer
optimization: better routing, improved tool selection priors, and refined
tier upgrade thresholds.

\subsection{Formal Objective}

The cross-layer optimization objective is:
\begin{equation}
    \max_{\bm{\theta}} \mathcal{J}(\bm{\theta})
    = \max_{\bm{\theta}} \mathbb{E}_{\tau \sim \mathcal{D}}\left[
        \text{TaskSuccess}(\tau; \bm{\theta}) - \lambda \cdot \text{Cost}(\tau; \bm{\theta})
    \right],
    \label{eq:joint-objective}
\end{equation}
where $\mathcal{D}$ is the empirical task distribution, $\lambda > 0$ is the
cost penalty coefficient, and $\text{TaskSuccess} \in [0, 1]$ is a graded
success measure. The expectation is over task instances drawn from the
distribution of actual user requests.

The UAH optimizes Equation~\ref{eq:joint-objective} via a combination of
gradient-free methods (for discrete routing decisions), structured prediction
(for tier selection), and Bayesian optimization (for the cost penalty
coefficient $\lambda$, which is tuned per user/organization based on their
stated cost-performance tradeoff).

% ============================================================================
% 6. COMPETITIVE ANALYSIS
% ============================================================================
